\input format.tex

\usepackage{graphicx}
\graphicspath{{cores/}}

\usepackage{environ}
\usepackage{colortbl,array,booktabs}
\usepackage{tabularx}
\usepackage{arydshln}
\newcolumntype{K}[1]{>{\centering\arraybackslash}p{#1}}

\usepackage[nodisplayskipstretch]{setspace}

\colorlet{TablaBordeSuperior}{topcolor}
\colorlet{TablaBordeInferior}{topcolor}
\colorlet{TablaCentroSuperior}{blue!1}
\colorlet{TablaCentroInferior}{blue!20}
\colorlet{FuenteCabeceraTabla}{white}

\newcolumntype{M}[1]{>{\centering\arraybackslash}m{#1}}

\tcbset{rtab/.style={
freelance,
frame code={
 \path[top color=topcolor,bottom color=topcolor]
   ([yshift=-#1*(\baselineskip+2pt)]interior.north west) --
   ([yshift=-#1*(\baselineskip+2pt)]interior.north east) {[sharp corners]--
    ([yshift=3pt]interior.north east) --
    ([yshift=3pt]interior.north west)} -- cycle;

  },
interior code={},
 }
}

\newcommand\fuentecabecera[1]{\textcolor{black}{\textbf{#1}}}


\begin{document}

\vspace*{5mm}
%% 各章节
\setlength{\arrayrulewidth}{1pt}
\fontsize{9.3pt}{11pt}\selectfont
\color{gray2}

\begin{tctabularx}{tabularx={m{15.4cm}<{\centering}}}
\\[-6pt]
\cellcolor{topcolor} \raisebox{2.614pt}{\color{black70}\fontsize{11pt}{11pt}\selectfont\bf 健康建议}
\end{tctabularx}

\sepline

\vspace*{1mm}
\begin{spacing}{1.5}
\noindent
\tikz\draw[topcolor,fill=topcolor](0,0) rectangle(0.1,0.25);
{\fontsize{10pt}{11pt}\selectfont\bf {综合您本次肠道有益菌检测，您肠道健康保护力总得分为14.27分，
处于正常范围，
建议您：}}
\end{spacing}

\vspace*{-4mm}
\noindent
\fontsize{8pt}{11pt}\selectfont
\arrayrulecolor{gray2}
\setlength{\arrayrulewidth}{.5pt}
\begin{center}
\begin{tabular}{|K{2.92cm}K{12.0cm}|}
\hline
\parbox[c][4.7cm]{.95\hsize}{
\noindent
\includegraphics[scale=1.0]{baohulifoodwithvline.pdf}
}
 &
\hspace*{4mm}
\parbox{.95\hsize}{
\vspace*{3mm}
\begin{spacing}{1.5}
\includegraphics[scale=1.0]{xiaofangkuai.pdf}
{\fontsize{8pt}{11pt}\selectfont 保持良好的饮食结构。建议保持定时定量进餐的好习惯，早餐不可忽略。食物种类新鲜多样，交替食用。主食搭配适量粗粮，多食用各种新鲜的蔬菜水果，深绿色蔬菜至少占蔬菜总摄入量的一半。控制油炸类、西式快餐、含糖饮料等高脂高糖高热量食物。忌暴饮暴食。减少外出就餐的次数。戒烟限酒。}
\end{spacing}
} \\
\hline

\parbox[c][4.7cm]{.95\hsize}{
\noindent
\includegraphics[scale=1.0]{baohulipreprobioticswithvline.pdf}
}
 &
\hspace*{4mm}
\parbox{.95\hsize}{
\vspace*{3mm}
\begin{spacing}{1.5}
\includegraphics[scale=1.0]{xiaofangkuai.pdf}
{\fontsize{8pt}{11pt}\selectfont 恭喜您，您的肠道保护力正常，建议您平时适量多喝酸奶。酸奶营养丰富易吸收，含有大量活性益生菌的酸奶，能帮助调节肠道酸碱值，抑制有害菌的生长，增强人体免疫力。选酸奶有窍门，2型糖尿病、动脉粥样硬化、胆囊炎、胰腺炎等患者，宜选用无糖酸奶。}
\end{spacing}
} \\
\hline
\parbox[c][4.7cm]{.95\hsize}{
\noindent
\includegraphics[scale=1.0]{baohuliotherswithvline.pdf}
}
 &
\hspace*{4mm}
\parbox{.95\hsize}{
\vspace*{3mm}
\begin{spacing}{1.5}
\includegraphics[scale=1.0]{xiaofangkuai.pdf}
{\fontsize{8pt}{11pt}\selectfont 建议戒烟限酒；加强体育运动，保持正常体重；正常作息，避免熬夜；同时，注意释放工作和生活压力，避免过重心理负担。}
\\
\includegraphics[scale=1.0]{xiaofangkuai.pdf}
{\fontsize{8pt}{11pt}\selectfont 建议您进行更详细的肠道菌群检测并持续监测，获取更具有针对性的膳食方案、微生态制剂推荐及运动指导等健康管理建议，为您的肠道健康保驾护航。}
\end{spacing}
} \\
\hline
\end{tabular}
\end{center}

{\noindent\qihao *温馨提示：以上健康建议仅供参考，您的实际膳食、运动及营养补充需根据您的自身具体情况进行调整。}
\vspace*{1cm}

\end{document}
